\section{Pathway-Level Analysis}

§3–§7 按转化路径逐一梳理了社区将 agent experience 从一种载体重新表示为另一种载体的机制。本章横向比较这些路径与它们产出的载体：每条 pathway 在生产侧付出什么，每类 carrier 在利用侧表现如何，任务约束如何在二者之间做出选择。

Section~\ref{sec:agent} 将 agent 决策写为 $a_t \sim \pi_\theta(\cdot \mid c_t)$，经验影响决策只有两条途径——流经 $c_t$ 的外部信息，或固化在 $\theta$ 的内部能力。Section~\ref{sec:carriers} 据此将载体分为 Tokenized， Latent 和 Parametric：输入侧的离散 token、输入侧的连续向量、权重内的隐式表示。

Section~\ref{sec:prod_profiles} 以五个维度刻画每条 pathway 的生产画像，Section~\ref{sec:util_profiles} 以四个维度刻画每类 carrier 的利用画像，Section~\ref{sec:selection_constraints} 将任务约束读为这两组维度的对偶——分析约束如何剪除不可行载体、约束冲突时如何迫使载体组合、约束随部署漂移时最优配置如何随之改变。

\subsection{Pathway Production Profiles}
\label{sec:prod_profiles}

本节以五个维度刻画每条 pathway 的生产画像。Production Mechanism（D1）：转化靠什么计算机制执行，如 LLM inference、gradient-based training、gradient-free persistence、execution + verification loop。Construction Cost（D2）：转化消耗多少资源。Information Transformation（D3）：经验信息在转化中被保留、损失、注入了什么，以及能否被追溯（traceability）；其中 loss 区分 lossy-by-design（抽象/压缩必然丢弃细节，损失类别可控）与 lossy-by-limitation（欠拟合、reward hacking、采样偏差，损失类别不可控，属工程风险而非设计意图）。Incrementality（D4）：新经验持续到来时，转化能否增量执行。Output Verifiability（D5）：转化产物能否在不依赖下游任务表现的情况下被独立验证。





\begin{landscape}

{\small                          % 字号调小，适配宽表
\setlength{\tabcolsep}{5pt}      % 列间距微调

\begin{longtable}{%
  L{2.0cm}   % Pathway
  L{2.8cm}   % Mechanism
  L{2.5cm}   % Construction Cost
  L{6.5cm}   % Information Transformation
  L{3.5cm}   % Incrementality
  L{2.8cm}   % Output Verifiability
  L{2.2cm}   % 代表工作 / 详节
}

%── 标题（如处于第 8 章，LaTeX 会自动编号为 Table 8.1）──────
\caption{Pathway production profiles across five dimensions.}%
\label{tab:pathway-profiles}\\

%── 首页页眉 ──────────────────────────────────────────────
\toprule
\textbf{Pathway} &
\textbf{Mechanism} &
\textbf{Construction Cost} &
\textbf{Information Transformation} &
\textbf{Incrementality} &
\textbf{Output Verifiability} &
\textbf{代表工作 / 详节} \\
\midrule
\endfirsthead

%── 续页页眉 ──────────────────────────────────────────────
\multicolumn{7}{c}{\tablename~\thetable\quad（续表）}\\[3pt]
\toprule
\textbf{Pathway} &
\textbf{Mechanism} &
\textbf{Construction Cost} &
\textbf{Information Transformation} &
\textbf{Incrementality} &
\textbf{Output Verifiability} &
\textbf{代表工作 / 详节} \\
\midrule
\endhead

%── 续页页脚 ──────────────────────────────────────────────
\midrule
\multicolumn{7}{r}{\small 接下页} \\
\endfoot

%── 末页页脚 ──────────────────────────────────────────────
\bottomrule
\endlastfoot

%── P1 ───────────────────────────────────────────────────
P1 Narrative Abstraction &
LLM inference &
低·可全自动 &
\textbf{保留}: 模式与洞察 /
\textbf{损失}: 步骤与时序(by-design) /
\textbf{注入}: 跨轨迹综合 insight(源自模型先验) &
增量(append); dynamic store 有 merge/prune 级联 &
直接阅读判断(主观) &
Reflexion [Shi23b]; \S3.1 \\
\midrule
%── P2 ───────────────────────────────────────────────────
P2 Schematic Formalization &
LLM inference + 可选 execution 验证 &
中·需可执行环境 &
\textbf{保留}: 过程结构 /
\textbf{损失}: NL nuance/隐含假设(by-design) /
\textbf{注入}: 显式依赖结构 &
增量(新 artifact 入库) &
execution test(客观, coverage-limited) &
Voyager [Wan23c]; \S3.2 \\
\midrule
%── P3 ───────────────────────────────────────────────────
P3 Latent-Space Transformation &
cache: gradient-free / trained: gradient &
cache 低·仅 session; trained 中 &
\textbf{保留}: 计算状态 /
\textbf{损失}: 符号可读性(by-design) /
\textbf{注入}: token 未显式编码的统计规律 &
cache 即时重建; trained 需 batch 重训 &
无直接验证·仅 ablation 间接 &
LAG [Che25b](cache)/
ReasonCACHE [Gup26](trained); \S4 \\
\midrule
%── P4 ───────────────────────────────────────────────────
P4 Evaluator Internalization &
gradient training &
高·需标注或明确成败信号 &
\textbf{保留}: 评估偏好 /
\textbf{损失}: 校准来源与单条贡献
(mix: by-design + by-limitation) &
batch 重训; 新标注需重训以保校准 &
无直接验证·hold-out calibration check &
Let's Verify Step by Step [Lig23]; \S5.1 \\
\midrule
%── P5 ───────────────────────────────────────────────────
P5 Policy Internalization &
gradient training (SFT/RL) &
高·需大规模高质量 trajectory &
\textbf{保留}: 决策模式 /
\textbf{损失}: 单条贡献不可区分(by-design)
+ 欠拟合/reward hacking 风险(by-limitation) &
batch 重训(可热启动); RL 可在线但漂移衰减旧经验 &
仅 behavioral (rollout 间接比较) &
SWE-Gym [Pan24]; \S5.2 \\
\midrule
%── P6 ───────────────────────────────────────────────────
P6 Evaluator-Driven Optimization &
gradient training (RLHF/DPO) &
高·叠加 P4 成本 &
\textbf{保留}: 偏好迁移入 policy /
\textbf{注入}: trajectory-only 学不到的细粒度偏好 /
\textbf{损失}: 经 P4+P6 双重压缩 &
online/iterative·漂移敏感 &
behavioral·循环验证风险
(Evaluator $\rightleftarrows$ Policy) &
Constitutional AI [Bai22b]; \S6.1 \\
\midrule
%── P7 ───────────────────────────────────────────────────
P7 Parametric Externalization &
从 policy 采样生成 &
中·依赖采样质量·须先有 trained policy &
\textbf{保留}: 权重中的能力经采样部分恢复为 token /
\textbf{损失}: 采样未覆盖部分 /
\textbf{注入}: hallucination 风险 &
可增量生成新样本·跟随源模型更新 &
产物可读但不可追溯权重成分·完备性不可验证 &
Explorer [Pah25]; \S6.2 \\

\end{longtable}

} % end \small

\end{landscape}

\textbf{Production Mechanism} 把七条路径沿是否更新权重切成两组。P1、P2、P3 的 cache 子类与 P7 不更新权重，其中 P1、P2 靠 LLM inference，P3-cache 靠 gradient-free 的前向状态留存，P7 靠从已训练模型采样；P3 的 trained 子类与 P4、P5、P6 经 gradient training，需数据筹备与优化的整套机制，单步资源消耗高于一次 inference。P2 在 inference 之外可叠加 execution + verification loop，是不更新权重一侧唯一带客观产出闸门的机制。

\textbf{Construction Cost} 在 gradient-based training 一类整体高、LLM inference 一类整体低。P4 卡在标注或明确成败信号的获取，P5 卡在大规模高质量 trajectory，P6 在 P4 的 Evaluator 训练之上再叠一层、成本累加。P1 可全自动、低成本，P3-cache 获取成本极低，P2 中等、需可执行环境；P7 受制于采样质量，且须先有 trained policy——它的低成本只就外化这一步而言，不含上游训练。

\textbf{Information Transformation} 的关键区分在 loss 的两种性质。七条路径注入的新信息都来自转化模型的先验、而非源经验。P1（全局的细节信息）、P2（丢 NL nuance）、P3（丢符号可读性）的损失均为 lossy-by-design，损失类别可控、由抽象策略调节；P4、P5 在 by-design 之外混入 lossy-by-limitation——P4 的标注噪声与标注者不一致直接进权重，P5 的 SFT 欠/过拟合与 RL reward hacking 不可控，属工程风险而非设计意图；P6 的损失是 P4 阶段与 P6 阶段的双重压缩，原始经验语义到达后已不可辨识；P7 的损失来自采样未覆盖的部分，并注入 hallucination 风险。关于 traceability：P3-cache 的 KV 直接来自真实前向传播、能对应具体历史片段，P3-trained 的 soft prompt 已难以对应回单条经验，P4/P5/P6 把单条贡献经 batch 梯度平均进权重、无法追问某个决策用了训练集中哪条轨迹；P7 产物虽是 Tokenized 端的离散 token，却采样自 Parametric 权重、追不到源知识。这是 batch 梯度平均的结构性后果，与工程精度无关。

\textbf{Incrementality} 在 inference 一类天然增量——P1、P2、P3-cache 产生的经验可支持增加，但随着规模的变大，有效利用难度上升。P7 随源模型更新持续生成新样本；gradient 一类须 batch 重训，P3-trained 改训练数据即重训整个 bank、P4 的新标注需重训以保校准。P5 可热启动、RL 还能在线更新，但旧经验的效用随分布漂移衰减，在线性以遗忘旧经验为代价；P6 的 online/iterative 更新同样对漂移敏感。P1 一侧的 dynamic store 在 append 之外带 merge/prune 级联，单条写入可能触发已有条目的连带改写。

\textbf{Output Verifiability} 取决于产物容许哪种独立检验。P2 的 code skill 由环境 execution 给出客观二值信号（如 Voyager 的 verification gate，通过/失败可判定），P1 的经验条目和P7的生成数据只能靠人或强 LLM 主观判断。Parametric 端，P4 的 Evaluator 尚可用 hold-out calibration 间接核验，P3、P5、P6 退到只能借 rollout 或 ablation 推断。

P6 与 P7 都以已固化的 Parametric 经验为源、做第二步加工，方向相反、在 pathway 图中互补：P6 把 Evaluator 的判断进一步写入 Policy 权重、深入 Parametric 端，P7 把 Policy 或 Evaluator 的能力采样回 Tokenized、退出 Parametric 端，使权重中的隐式经验重新成为可供其他路径使用的显式载体。循环验证与 hallucination 是仅现于转化过程的两类风险：前者来自 P6 中 Evaluator 与 Policy 互为参照、缺独立外部 ground truth（ECHO \cite{Li26l} 的同步更新可缓解、不能消除），后者来自 P7 从权重采样、缺环境 grounding 时产出似真而无据的 trajectory。

% Policy 由 P5、P6 两条路径产生，二者代价画像不同，使 Section~\ref{sec:selection_constraints} 对 Policy 目标的可行性判断成为 P5 还是 P6 的选择、而非二元判断；

\subsection{Carrier Utilization Profiles}
\label{sec:util_profiles}

利用画像刻画五类载体生产完成后如何进入决策循环。载体的进入方式由其存在形态决定，与生产它的路径无关：Latent 向量无论来自 cache 留存还是训练所得模块，都作为连续向量参与 forward pass。沿四个维度展开。Access Interface（D1）：经验从 $c_t$ 还是 $\theta$ 进入决策循环。Access Granularity（D2）：能否在推理时有选择地激活特定经验，还是整个载体始终全开。Utilization Cost（D3）：使用经验需支付什么资源、计算发生在何处。Revisability（D4）：经验有误时如何修正，修正是否波及其他经验。

% Schematic artifact 无论是 code、workflow 还是 graph，都是 $c_t$ 中的强形式化离散 token、靠结构被操作（parse、execution 或 traversal），其中 code-skill 形态把执行卸至外部 executor。

% 这些层内差异改变具体单元格的取值（利用代价、修正方式），但不改变载体的可寻址性与访问粒度，记为单元格内限定，不另立载体。

% 可解释性与控制自由度都是访问粒度的下游读数。可读性随粒度与层次走——Tokenized 因逐条可寻址加离散符号而可读，Parametric 因不可寻址加连续权重而不可读；控制自由度同样系于粒度——逐条可寻址允许在检索、编排、注入各环节调节，不可寻址则不允许。二者作为任务约束的对偶留待 §8.3。

% 所需宏包：
% \usepackage{booktabs}
% \usepackage{tabularx}
% \usepackage{ctex}        % 中文支持（XeLaTeX / LuaLaTeX）

\begin{table}[htbp]
  \centering
  \caption{Carrier utilization profiles across four dimensions.}
  \label{tab:8.2}
  \small
  \renewcommand{\arraystretch}{1.5}
  \begin{tabularx}{\textwidth}{lXXXX}
    \toprule
    \textbf{Carrier} &
    \textbf{Access Interface} &
    \textbf{Access Granularity} &
    \textbf{Utilization Cost} &
    \textbf{Revisability} \\
    \midrule
    Narrative &
    context injection（$c_t$） &
    per-item &
    Retrieval Cost + Inference Cost &
    - \\
    \midrule
    Schematic &
    context injection（$c_t$） &
    per-item &
    Executor Cost + Inference Cost &
    - \\
    \midrule
    Latent &
    forward-pass attention &
    per-item &
    Inference Cost &
    - \\
    \midrule
    Policy &
    forward-pass weights（$\theta$）$\rightarrow$ $a_t$ &
    always-on &
    Inference Cost &
    - \\
    \midrule
    Evaluator &
    forward-pass weights（$\theta$）$\rightarrow$ 候选信号 &
    always-on &
    Inference Cost &
    - \\
    \bottomrule
  \end{tabularx}
\end{table}


访问接口由存在形态锁定：Tokenized 走 context injection（$c_t$），Latent 走 forward-pass attention（仍在输入侧，但以连续向量参与），Parametric 走 forward-pass weights（$\theta$）。选定载体即锁定接口。访问粒度沿 Tokenized → Latent → Parametric 单调变粗：Tokenized 逐条取用（如取单条 reflection 或 skill），Latent 介于逐条和整体之间，Parametric 全局生效。

层内的两个分叉不改变这条单调性，各自只动一两个维度。Tokenized 内，Narrative 与 Schematic 同处 per-item、同样可读，差别仅在复用机制——Narrative 靠 language understanding 阅读，Schematic 靠 parse、execution 或 graph traversal 按结构操作；这一差别落在 Output Verifiability（§8.1）与利用延迟（executable Schematic 的执行同步阻塞），不动可寻址性。Parametric 内，Policy 与 Evaluator 共享同一份存储画像——都不可寻址、都 retrain 才能改、复用都零边际——差别纯在功能角色：Policy 的 forward pass 直接产 $a_t$，Evaluator 的产一个关于候选的信号，可被训练消费（→ Evaluator-Driven Optimization）或推理消费（对 $\pi_\theta$ 候选 rerank、Best-of-N 筛选、tree search 剪枝）。以 PRM \cite{Lig23} 为例，同一套 step-level score 既能作 RL 的 dense reward，又能作 Best-of-N 的 step 级筛选或 tree search 的剪枝准则——这种下游灵活是"评估型输出可被多处取用"的性质，选哪个 evaluator、信号注入到哪个环节、取 outcome 还是 step 粒度，是把功能模块接入决策环的编排（属 §7），不是 Evaluator 的存储比 Policy 更可寻址。

在复用代价这一维上，Tokenized 与 Parametric 两端的取值难以：Tokenized 的经验可以单条取用，但是每次需要重新载入 Context 进行计算，存在复用成本，且当经验库的规模较大时，存在检索成本。Latent 居中，以 Hidden States 的形式参与前向计算，但同时也存在取用成本，比如检索或生成特定的 Hidden States。Parametric 的经验全局生效，每次复用边际成本可认为为零。

\subsection{Constraint-Driven Selection}
\label{sec:selection_constraints}

经验的载体与转化路径由任务的约束界定。将任务约束沿5个轴展开：Build Cost （C1） 刻画转化经验带来的成本；Serving Cost（C2）刻画使用经验的代价；Inspectability（C3）刻画经验的可解释性，包含 readability 和 verifiability；Maintainability（C4）刻画经验的维护成本，包含错误经验的修正 （error correction）和环境的非稳定性变化（non-stationarity） ；Scale（C5）刻画经验的可扩展性；


\begin{table}[htbp]
  \centering
  \caption{Constraint axes, their sub-constraints, and the \S8.1/\S8.2 dimensions each engages.}
  \label{tab:8.3}
  \begin{tabularx}{\textwidth}{llXX}
    \toprule
    约束轴 & 子约束 & 利用侧维度（\S8.2） & 生产侧维度（\S8.1） \\
    \midrule
    Build Cost
      & training resources
      & ---
      & Construction Cost、Mechanism \\
    Serving Cost
      & latency / context budget
      & Utilization Cost、Access Interface
      & --- \\
    Inspectability
      & readability
      & Access Granularity
      & Information Transformation（traceability） \\
    Inspectability
      & verifiability
      & ---
      & Output Verifiability \\
    Maintainability
      & error correction / non-stationarity
      & Revisability
      & Incrementality \\
    Scale（放大轴）
      & volume / arrival rate
      & 放大各轴
      & 放大各轴 \\
    \bottomrule
  \end{tabularx}
\end{table}


\begin{table}[htbp]
  \centering
  \caption{Directional pressure from constraint sub-dimensions to carrier classes. $\mathtt{++}$\,强烈偏向，$\mathtt{+}$\,偏向，$\circ$\,中性，$\mathtt{-}$\,不利，$\mathtt{{-}{-}}$\,强烈不利,指载体在该约束下的适配度。}
  \label{tab:8.4}
  \begin{tabular}{ll*{5}{c}}
    \toprule
    约束轴 & 子约束 & Narrative & Schematic & Latent & Policy & Evaluator \\
    \midrule
    \multirow{2}{*}{Serving Cost}
      & latency
        & $-$ & $--$ & $\circ$ & $++$ & $++$ \\
    \cmidrule(l){2-7}
      & context budget
        & $-$ & $--$ & $+$ & $++$ & $++$ \\
    \midrule
    \multirow{2}{*}{Inspectability}
      & readability
        & $++$ & $++$ & $-$ & $--$ & $--$ \\
    \cmidrule(l){2-7}
      & verifiability
        & $++$ & $++$ & $-$ & $--$ & $--$ \\
    \midrule
    \multirow{2}{*}{Maintainability}
      & error correction
        & $++$ & $+$ & $-$ & $--$ & $--$ \\
    \cmidrule(l){2-7}
      & non-stationarity
        & $++$ & $+$ & $-$ & $--$ & $--$ \\
    \midrule
    Build Cost
      & training resources
        & $+$ & $++$ & $-$ & $--$ & $--$ \\
    \bottomrule
  \end{tabular}
\end{table}

\textbf{Build Cost} 由转化是否更新权重决定，资源含算力与标注两类。gradient-free 的转化——P1、P2 与 P3 的 cache 子类——预算受限时仍可执行，不付训练成本；gradient-based 的转化——P3 的 trained 子类与 P4、P5、P6——需完整训练循环，预算不足时不可行，Latent 在此沿 cache 与 trained 分化最明显。算力与标注可分别稀缺：纯环境可验证信号（unit test、ground-truth match）驱动的 Policy RL 不需人工标注，而依赖人类示范或偏好标签的 SFT 与 reward modeling 在标注稀缺时受限更重。

\textbf{Serving Cost} 分时间与空间两面，区别在该代价随每次复用重新产生、还是一次性摊销。Narrative 两面均付、且每次复用重付：条目须检索并 prefill 进 context \cite{Shi23b, Zha23c}。Schematic 分形态——executable 将计算卸至外部 executor、不占 context，但执行同步阻塞决策；non-executable 仍注入 context。Latent 较为特殊，最终的复用方式都是以向量的形态直接参与模型的前向计算，但是其复用成本即可通过检索已有记忆库选出，也可通过相关模块来临时生成。Policy 与 Evaluator 边际复用成本为零，权重常驻。

\textbf{Inspectability} 分为 readability 与 verifiability，前者要逐条可读，后者要使用前可验证。Narrative 逐条可读、可主观核验，两者俱强 \cite{Zha23c}。Schematic 分形态：non-executable 的描述性结构可读性高，可主管验证 \cite{Ano24}；executable 的程序需技术能力方能阅读、可读性相对较低，但其执行可在落库前客观测试 \cite{Wan23c}。Latent， Policy 与 Evaluator 既不可读也无法事前验证。

\textbf{Maintainability} error correction 是修正已知错误的速度，non-stationarity 是跟随漂移的能力。Tokenized 一般逐条独立，单条修正不触发全局重写，新增条目即时可检索 \cite{Zha23c,Wan23c}。 对错误的修正较快，且对环境变化的适应能力较好。Latent 两面俱弱于 Tokenized。 Parametric 的表现最差，修正需 retrain 并可能引发 catastrophic forgetting，其在线更新以近端经验覆盖旧经验、旧经验效用随之衰减 \cite{Bai24}。


\textbf{Scale} 经验体量放大其余各轴的压力。随着经验库的增大，Tokenized 的检索质量会下降 \cite{Fu24}，维护成本会非线性的增加，加入 Context 后的推理成本接近是固定的；Parametric 推理成本与训练数据体量解耦、训练完成后保持恒定 \cite{Pan24}，其构建成本与维护成本会线性增加。 Latent 既存在检索使用的方式，有存在需要由训练好的模型临时生成的方式，介于两者之间。


约束为对应维度设定阈值，载体能否进入决策循环，取决于其取值是否满足。约束相容时逐一剪枝即可定位可行载体；约束冲突时单载体无解，迫使载体在时间或空间上组合。约束还会被误判、随部署漂移，使最优配置成为时间的函数。